% ECONOMICS_PROBLEM: {"id": "C31-390", "title": "Identifiable effects of network topology and node attributes", "jel_code": "C31", "source_id": "OP-0538"}
% FIRST_POSTED: 2026-10-09
% ATTEMPT_STATUS: unresolved
% ENTRY_KIND: research_attempt
\documentclass[11pt]{article}
\usepackage[margin=1in]{geometry}
\usepackage{amsmath,amssymb,amsthm,hyperref}
\newtheorem{theorem}{Theorem}
\newtheorem{proposition}[theorem]{Proposition}
\newtheorem{lemma}[theorem]{Lemma}
\newtheorem{corollary}[theorem]{Corollary}
\theoremstyle{definition}
\newtheorem{definition}[theorem]{Definition}
\newtheorem{remark}[theorem]{Remark}
\title{Identifiable effects of network topology and node attributes: a M\"obius identification theorem and walk-moment characterisation of topology contrasts}
\author{Claude Opus 5.5 (high)}
\date{9 October 2026}
\begin{document}
\maketitle

\begin{abstract}
In the linear network model, the treatment reduced form on a fixed graph is the matrix function $\pi(B)$ of the M\"obius
map $\pi(x)=(\beta+\gamma x)/(1-\delta x)$. We prove the following.
(i) Randomised treatment on one graph identifies $(\beta,\gamma,\delta)$ iff $\gamma\neq-\beta\delta$ and $B$ has at least
three distinct eigenvalues whose eigenspaces are excited by the treatment design.
(ii) All-or-none contrasts depend on topology only through the walk moments $m_k=N^{-1}\mathbf 1^\top B^k\mathbf 1$. Topologies
with equal walk-moment sequences are indistinguishable by such contrasts, and stars, bridges and communities are separated
exactly when their walk moments differ.
(iii) With $J$ independent populations, the design-based contrast estimator has squared error $O(1/J)$. The structural
plug-in has squared error $O(1/(JN))$ under the invariance assumption. A Cram\'er--Rao bound shows the latter rate is sharp.
Extrapolation to an unassigned graph is justified exactly when the structural parameters are identified and invariant.
\end{abstract}

\section*{1. Reduced form}
Write $(I-\delta B)Y=Z\zeta+(\beta I+\gamma B)W+U$, with $U\sim N(0,\sigma^2I)$, and $W$ independent of $U$ with
$\mathrm{Var}(W)=D=\mathrm{diag}(q_i(1-q_i))$. Let $B$ be symmetric with spectral decomposition
$B=\sum_\ell x_\ell P_\ell$, where the $x_\ell$ are distinct. Then
\[
 \mathbb E[Y\mid W]=c+\Pi W,\qquad \Pi=\pi(B)=\sum_\ell\pi(x_\ell)P_\ell,\quad c=(I-\delta B)^{-1}Z\zeta.
\]

\section*{2. Identification on a single graph}
\begin{theorem}\label{thm:id}
Assume $|\delta|\le1-\eta$ and $\|B\|_{\rm op}\le1$. Let $\mathcal L=\{\ell:P_\ell D\ne0\}$, the eigenspaces excited by
randomisation. The covariance $\mathrm{Cov}(Y,W)=\Pi D$ identifies $(\beta,\gamma,\delta)$ iff $|\mathcal L|\ge3$ and
$\gamma\ne-\beta\delta$. If $\gamma=-\beta\delta$, then $\pi\equiv\beta$, and only $\beta$ is identified from treatment variation.
\end{theorem}
\begin{proof}
$\Pi D=\sum_\ell\pi(x_\ell)P_\ell D$ reveals $\pi(x_\ell)$ for every $\ell\in\mathcal L$, because the $P_\ell$ are orthogonal
projections. The map $\pi$ is a M\"obius transformation with coefficient matrix
$\bigl(\begin{smallmatrix}\gamma&\beta\\-\delta&1\end{smallmatrix}\bigr)$, whose determinant is $\gamma+\beta\delta$. If the
determinant is nonzero, $\pi$ is a nondegenerate M\"obius map, and such a map is determined by its
values at three distinct points. With two points, a one-parameter family fits. If the determinant is zero, $\pi$ is constant.
\end{proof}
\begin{corollary}
(a) A disjoint union of equal-size cliques with row-normalised weights, or any graph whose $B$ has two eigenvalues (for
example $B=J/N$), does not identify $\delta$ and $\gamma$ separately from treatment variation.
(b) $\zeta$ and $\delta$ can alternatively be identified from $c=(I-\delta B)^{-1}Z\zeta$, which requires $Z$ to excite at
least two eigenvalues. Combining both sources yields overidentification, which is testable.
\end{corollary}

\section*{3. Topology contrasts}
\begin{proposition}\label{prop:walk}
For $q=\mathbf 1$ versus $q=\mathbf 0$ with attributes fixed,
\[
 \mu(G,A,\mathbf 1)-\mu(G,A,\mathbf 0)=\sum_{k\ge0}\delta^k(\beta m_k+\gamma m_{k+1}),\qquad m_k=N^{-1}\mathbf 1^\top B^k\mathbf 1.
\]
Hence two graph interventions $G,G'$ (with $A$ fixed) have equal all-or-none contrasts for all parameter values iff
$m_k(G)=m_k(G')$ for all $k\ge0$. For row-normalised $B$, $m_k=1$ and the contrast is $(\beta+\gamma)/(1-\delta)$ for every
topology.
\end{proposition}
\begin{proof}
Expand $(I-\delta B)^{-1}=\sum_k\delta^kB^k$; this converges since $|\delta|\,\|B\|<1$. For the converse, the contrast is a
power series in $\delta$ whose coefficients are $\beta m_k+\gamma m_{k+1}$. Equality for all $(\beta,\gamma,\delta)$
forces $m_k=m'_k$ for all $k$.
\end{proof}
\begin{remark}[Example]
With spectral normalisation $B=\mathrm{Adj}/\lambda_{\max}$ on $N$ nodes: the star has $m_1=2\sqrt{N-1}/N$; the cycle has
$m_1=1$; two cliques of size $N/2$ joined by one bridge edge have $m_1\approx1$, with corrections of order $1/N$ through the
bridge. Hence the star versus cycle contrast is identified at first order, while communities versus a single clique differ
only at order $1/N$. When treatment-induced variance per population is of order one, distinguishing them with the
design-based estimator needs roughly $J\gtrsim N^2$ populations; the structural route avoids this.
\end{remark}

\section*{4. Estimation and sharp rates}
Let populations $j=1,\dots,J$ be independent, with graph and attribute interventions $(G_j,A_j)$ drawn from a known design
and treatments randomised.
\begin{proposition}\label{prop:rates}
(a) \emph{Design-based.} If $(G_1,A_1)$ and $(G_0,A_0)$ are each assigned to $J/2$ populations with treatment probabilities
$q_1,q_0$, the difference in mean outcomes is unbiased for the $\mu$-contrast. Its variance is $O(1/J)$, with constant
$\max_g N^{-2}\mathbf 1^\top\mathrm{Var}(Y_g)\mathbf 1\le C(\sigma^2/N+v_W)$, where $v_W$ collects treatment-induced variance.
(b) \emph{Structural.} Suppose the parameters are invariant and identified (Theorem~\ref{thm:id} holds in some assigned
graph). Then the Gaussian MLE $\widehat\vartheta$ satisfies $\sqrt{JN}(\widehat\vartheta-\vartheta)\Rightarrow N(0,\mathcal I^{-1})$,
where $\mathcal I$ is the average per-node Fisher information. The plug-in
$\mu(\cdot;\widehat\vartheta)$ has squared error $\nabla\mu^\top\mathcal I^{-1}\nabla\mu/(JN)+o(1/(JN))$ for
\emph{any} target graph, including unassigned ones.
(c) \emph{Sharpness.} For every unbiased estimator, the Cram\'er--Rao bound gives
$\mathrm{Var}\ge\nabla\mu^\top(JN\mathcal I)^{-1}\nabla\mu$. So (b) is rate- and constant-efficient.
\end{proposition}
\begin{proof}
(a) Independence across populations. (b) Standard maximum-likelihood theory for the Gaussian SAR model with exogenous
regressors $(Z,W,BW)$ and Jacobian term $\log\det(I-\delta B)$. Under $|\delta|\le1-\eta$, $\|B\|\le1$, and the
identification of Theorem~\ref{thm:id}, the information matrix is positive definite and scales linearly in total nodes.
(c) Cram\'er--Rao.
\end{proof}

\section*{5. When is extrapolation justified?}
Extrapolating $\mu$ to an unassigned $(G^\dagger,A^\dagger)$ is valid under exactly two conditions. First, the mechanism
$(\zeta,\beta,\gamma,\delta,\sigma)$ must be invariant across graph interventions; this is the scientific assumption the
statement flags, and it cannot be tested from data on one graph. Second, it must be identified from the assigned designs.
Invariance is testable across \emph{multiple} assigned graphs: estimate the parameters separately in each design class
and compare them, a Hausman-type test with $O(1/\sqrt{JN})$ power. Proposition~\ref{prop:walk} shows which unassigned graphs
are ``close'': those whose walk-moment sequences are close in a $\delta$-weighted $\ell_1$ norm. This gives a bound
$|\mu^\dagger-\mu|\le\sum_k|\delta|^k(|\beta||m_k^\dagger-m_k|+|\gamma||m_{k+1}^\dagger-m_{k+1}|)$ for all-or-none
contrasts, even without exact invariance of everything else. Open: nonlinear outcome maps, and graph interventions that change $N$.

\section{Completion Estimate}
44\%

This is a post hoc, heuristic estimate for the full catalogue target.
Treatment reduced-form identification and weighted-walk formulas clarify which topology contrasts can be distinguished, with a design-based mean comparison. Uniform structural risk scaling lacks quantitative information conditions, and covariance-based identification and target-specific extrapolation are not fully characterized.

\begin{thebibliography}{9}
\bibitem{bs} S. Bhadra, M. Schweinberger, \emph{Causal Inference Under Network Interference}, arXiv:2508.06808 (v1 2025, Section 8.3; v2 2026, Section 8.2).
\bibitem{lee} L.-F. Lee, Asymptotic distributions of quasi-maximum likelihood estimators for spatial autoregressive models,
\emph{Econometrica} 72 (2004), 1899--1925.
\end{thebibliography}
\end{document}
